\section{Open Questions}

Five genuinely-open, load-bearing questions from re-reading the paper. Each is grounded in a specific gap between what the paper claims and what independent verification can currently support.

\subsection*{Q1. What MCDS input $Y(\text{LET})$ table did Wang actually use?}

\textbf{Basis.} The 6 fit parameters + published D10 only self-consistently reproduce when $Y_X$ is treated as a free $\approx$55 DSB/Gy/cell (HSG) or $\approx$51 (V79). McMahon 2017 MCDS gives $\approx$34, and Kukla's fresh MCDS 3.10A run gives $\approx$8 at low LET (all order-of-magnitude smaller). No independent MCDS calibration reproduces the paper's D10 --- the $\approx$55 value is unexplained by any published calibration and is nowhere printed in the paper. \emph{This is the load-bearing gap that keeps the verdict PARTIAL.}

\textbf{Next steps.} (1) Email corresponding author (Junli Li, Tsinghua) to request the MCDS 3.x configuration file and $Y(\text{LET}), \lambda(\text{LET})$ tables used for Table 1 and Figs.\ 2--5. (2) Install MCDS 3.10A, replicate Wang's nucleus geometry (HSG 6 Gbp, V79 5.6 Gbp) with matching cluster definitions, and generate $Y(\text{LET})$ for all 5+ ion species. (3) Refit the 6 parameters on the 106 PIDE curves with this independent $Y(\text{LET})$ --- if the resulting Table 1 differs but achieves similar in-sample $R^2$, the ambiguity is resolved and the model is confirmed under-identified.

\subsection*{Q2. Does the model's RBE prediction extrapolate to LETs outside the Furusawa fit range?}

\textbf{Basis.} The paper fits and evaluates on the same 106 Furusawa curves (in-sample $R^2$ only). Fig.\ 5's mixed-beam claim ($R^2=0.762$) is the closest to a generalization test, but mixed beams are still combinations of the base spectra used in fitting. True out-of-sample RBE at unfit LETs --- e.g.\ very low-LET proton ($<0.5$ keV/$\mu$m), very high-LET Fe-56 ($>1000$ keV/$\mu$m), or low-energy alpha at $\sim$150 keV/$\mu$m --- is not reported. Given the model's phenomenological character, extrapolation confidence is unclear.

\textbf{Next steps.} (1) Leave-one-out cross-validation over the 5 ions in PIDE, measuring predictive $R^2$ on the held-out set. (2) Compare against non-Furusawa datasets: Suzuki 2000 (alpha on HSG), Belli 2000 (proton on V79), Weyrather 1999 (C-12 on Chinese hamster). (3) Report cross-dataset RBE prediction $R^2$ --- a meaningful out-of-sample test the paper never performs.

\subsection*{Q3. Would DSB complexity outperform DSB number as predictor?}

\textbf{Basis.} The model uses only $Y$ (DSB/Gy) and $\lambda$ (DSBs/track), collapsing all cluster-complexity information to a single scalar per track. MCDS 3.x explicitly outputs cluster-size distributions $P(k)$ which are known from Nikjoo/Goodhead work to matter for lethality: 2+ break clusters are $\sim 4\times$ harder to repair than isolated DSBs. Wang's model absorbs this into the phenomenological $\zeta,\xi$ parameters but never tests whether an explicit complexity term improves fit.

\textbf{Next steps.} (1) Extend MCDS output to include full $P(k \mid \text{LET, ion})$ for the 5 ion species. (2) Extend the model with an explicit lethality-per-complexity term $N_{\text{death,complex}} = \sum_k w_k f_k Y D$. (3) Refit on PIDE 106 curves; measure whether adding this term improves in-sample $R^2$ and out-of-sample RBE prediction. Report whether Wang's 6-parameter fit is at the noise floor or leaving explainable variance on the table.

\subsection*{Q4. Can the model extend to hypoxia via OER-scaled $Y$?}

\textbf{Basis.} The paper addresses only aerobic clonogenic survival. Clinically the OER (oxygen enhancement ratio) is dose-modifying and LET-dependent (decreases from $\sim$3 at low LET to $\sim$1 at very high LET). A mechanistic DSB-based model \emph{should} naturally predict this via reduced $Y$ under hypoxia. Whether Wang's fit parameters transfer unchanged under a hypoxia-scaled $Y$, or need refitting, is unknown.

\textbf{Next steps.} (1) Use MCDS's built-in oxygen-fraction parameter to generate $Y_{\text{hypoxic}}(\text{LET})$ for pO$_2\sim 0.5\%, 2\%, 5\%, 21\%$. (2) Predict survival curves at each pO$_2$ with Table 1 parameters unchanged (only $Y$ changes). (3) Compare against hypoxic-survival datasets (Furusawa 2000 hypoxic conditions; Wouters 1997; Nagasawa 2003). (4) Report whether the aerobic-only fit generalizes --- a significant validation and extension if yes; a diagnostic if no.

\subsection*{Q5. Does full track-structure DSB spectrum (vs.\ collapsed $\lambda$) matter at clinical proton LETs ($\sim$2--5 keV/$\mu$m)?}

\textbf{Basis.} Proton-therapy RBE variation across the Bragg peak is a live clinical controversy (constant-RBE=1.1 vs.\ variable-RBE models). Wang uses $\lambda$ = DSBs/track as a single scalar per LET, but MCDS produces a full distribution $P(m \text{ DSBs} \mid \text{one track at LET }L)$. Collapsing to a mean $\lambda$ may lose predictive power exactly where clinical stakes are highest (proton distal edge, $\sim$5--10 keV/$\mu$m). Wang's Fig.\ 4 shows RBE mostly in the ion-therapy LET regime, not the sub-10 keV/$\mu$m proton regime.

\textbf{Next steps.} (1) Modify the model to consume the full $P(m \mid \text{LET})$ distribution rather than mean $\lambda$, keeping the Poisson misrepair framework. (2) Compare predictions against proton-specific benchmarks (McNamara 2015, Wedenberg 2013, R{\o}rvik 2018). (3) Compute predicted RBE at 2, 3, 5, 8 keV/$\mu$m and compare to the Paganetti clinical compilation. (4) Report whether track-spectrum-aware Wang matches or improves on the proton-RBE benchmarks --- the highest-clinical-value extension.
